\SecondArticle{16}
\textsc{Satz.} \textit{Es bezeichne $\varphi(u,x)$}\HistoricalNote{12}\textit{ ein Produkt aus beliebig vielen Faktoren, in welche die Unbestimmten $u,x$ nur linear eingehen, die also die Form}
\[
\begin{gathered}
\alpha+\beta u+\gamma x\\
\alpha'+\beta'u+\gamma'x\\
\alpha''+\beta''u+\gamma''x\\
\text{usw.}
\end{gathered}
\]
\textit{haben; ferner sei $w$ eine weitere Unbestimmte. Dann ist die Funktion}
\[
\varphi(u+w\, .\,\tfrac{\mathrm d\varphi(u,x)}{\mathrm dx},\ x-w\, .\,\tfrac{\mathrm d\varphi(u,x)}{\mathrm du})=\Omega
\]
\textit{identisch durch $\varphi(u,x)$ teilbar.}



\SourceLeaf{059}
\textit{Beweis.} Setzen wir
\[
\begin{aligned}
\varphi(u,x)&=(\alpha+\beta u+\gamma x)Q\\
&=(\alpha'+\beta'u+\gamma'x)Q'\\
&=(\alpha''+\beta''u+\gamma''x)Q''\\
&\quad\text{usw.}
\end{aligned}
\]
so sind $Q,Q',Q''$ usw. Polynome in den Unbestimmten $u,\allowbreak x,\allowbreak \alpha,\allowbreak \beta,\allowbreak \gamma,\allowbreak \alpha',\allowbreak \beta',\allowbreak \gamma'$, $\alpha'',\beta'',\gamma''$ usw., und es gilt
\[
\begin{aligned}
\frac{\mathrm d\varphi(u,x)}{\mathrm dx}&=\gamma Q+(\alpha+\beta u+\gamma x)\, .\,\frac{\mathrm dQ}{\mathrm dx}\\
&=\gamma'Q'+(\alpha'+\beta'u+\gamma'x)\, .\,\frac{\mathrm dQ'}{\mathrm dx}\\
&=\gamma''Q''+(\alpha''+\beta''u+\gamma''x)\, .\,\frac{\mathrm dQ''}{\mathrm dx}\\
&\quad\text{usw.}
\end{aligned}
\]
\[
\begin{aligned}
\frac{\mathrm d\varphi(u,x)}{\mathrm du}&=\beta Q+(\alpha+\beta u+\gamma x)\, .\,\frac{\mathrm dQ}{\mathrm du}\\
&=\beta'Q'+(\alpha'+\beta'u+\gamma'x)\, .\,\frac{\mathrm dQ'}{\mathrm du}\\
&=\beta''Q''+(\alpha''+\beta''u+\gamma''x)\, .\,\frac{\mathrm dQ''}{\mathrm du}\\
&\quad\text{usw.}
\end{aligned}
\]
Setzt man diese Werte in die Faktoren ein, deren Produkt $\Omega$ ist, nämlich in
\[
\begin{gathered}
\alpha+\beta u+\gamma x+\beta w\, .\,\frac{\mathrm d\varphi(u,x)}{\mathrm dx}-\gamma w\, .\,\frac{\mathrm d\varphi(u,x)}{\mathrm du}\\
\alpha'+\beta'u+\gamma'x+\beta'w\, .\,\frac{\mathrm d\varphi(u,x)}{\mathrm dx}-\gamma'w\, .\,\frac{\mathrm d\varphi(u,x)}{\mathrm du}\\
\alpha''+\beta''u+\gamma''x+\beta''w\, .\,\frac{\mathrm d\varphi(u,x)}{\mathrm dx}-\gamma''w\, .\,\frac{\mathrm d\varphi(u,x)}{\mathrm du}\\
\text{usw., jeweils}
\end{gathered}
\]
so erhalten diese die folgenden Werte:
\[
\begin{gathered}
(\alpha+\beta u+\gamma x)(1+\beta w\, .\,\tfrac{\mathrm dQ}{\mathrm dx}-\gamma w\, .\,\tfrac{\mathrm dQ}{\mathrm du})\\
(\alpha'+\beta'u+\gamma'x)(1+\beta'w\, .\,\tfrac{\mathrm dQ'}{\mathrm dx}-\gamma'w\, .\,\tfrac{\mathrm dQ'}{\mathrm du})\\
(\alpha''+\beta''u+\gamma''x)(1+\beta''w\, .\,\tfrac{\mathrm dQ''}{\mathrm dx}-\gamma''w\, .\,\tfrac{\mathrm dQ''}{\mathrm du})\\
\text{usw.}
\end{gathered}
\]
\SourceLeaf{060}
Daher ist $\Omega$ das Produkt aus $\varphi(u,x)$ und den Faktoren
\[
\begin{gathered}
1+\beta w\, .\,\frac{\mathrm dQ}{\mathrm dx}-\gamma w\, .\,\frac{\mathrm dQ}{\mathrm du}\\
1+\beta'w\, .\,\frac{\mathrm dQ'}{\mathrm dx}-\gamma'w\, .\,\frac{\mathrm dQ'}{\mathrm du}\\
1+\beta''w\, .\,\frac{\mathrm dQ''}{\mathrm dx}-\gamma''w\, .\,\frac{\mathrm dQ''}{\mathrm du}
\end{gathered}
\]
usw., also aus $\varphi(u,x)$ und einem Polynom in den Unbestimmten $u,\allowbreak x,\allowbreak w,\allowbreak\alpha,\allowbreak\beta,\allowbreak\gamma$, $\alpha',\beta',\gamma',\alpha'',\beta'',\gamma''$ usw. Q. E. D.

\Needspace{20\baselineskip}\SecondArticle{17}
Der Satz des vorigen Artikels ist offenbar auf die Funktion $\zeta$ anwendbar, die wir von nun an durch
\[
f(u,x,\lambda',\lambda'',\lambda'''\text{ usw.})
\]
darstellen. Dann ist
\[
f(u+w\, .\,\tfrac{\mathrm d\zeta}{\mathrm dx},\ x-w\, .\,\tfrac{\mathrm d\zeta}{\mathrm du},\ \lambda',\lambda'',\lambda'''\text{ usw.})
\]
identisch durch $\zeta$ teilbar. Den Quotienten, ein Polynom in den Unbestimmten $u,x,w,a,b,c$ usw., das hinsichtlich $a,b,c$ usw. symmetrisch ist, stellen wir durch
\[
\psi(u,x,w,\lambda',\lambda'',\lambda'''\text{ usw.})
\]
dar. Daraus schließen wir, dass auch identisch
\[
\begin{aligned}
&f(u+w\, .\,\tfrac{\mathrm dz}{\mathrm dx},\ x-w\, .\,\tfrac{\mathrm dz}{\mathrm du},\ l',l'',l'''\text{ usw.})\\
&=z\psi(u,x,w,l',l'',l'''\text{ usw.})
\end{aligned}
\]
sowie
\[
\begin{aligned}
&f(u+w\, .\,\tfrac{\mathrm dZ}{\mathrm dx},\ x-w\, .\,\tfrac{\mathrm dZ}{\mathrm du},\ L',L'',L'''\text{ usw.})\\
&=Z\psi(u,x,w,L',L'',L'''\text{ usw.})
\end{aligned}
\]
gilt. Stellen wir also die Funktion $Z$ einfach durch $F(u,x)$ dar, sodass
\[
f(u,x,L',L'',L'''\text{ usw.})=F(u,x)
\]
gilt, so haben wir identisch
\[
\begin{aligned}
&F(u+w\, .\,\tfrac{\mathrm dZ}{\mathrm dx},\ x-w\, .\,\tfrac{\mathrm dZ}{\mathrm du})\\
&=Z\psi(u,x,w,L',L'',L'''\text{ usw.})
\end{aligned}
\]

\SourceLeaf{061}
\SecondArticle{18}
Nehmen wir nun an, aus bestimmten Werten von $u,x$, nämlich aus $u=U,x=X$, ergebe sich
\[
\frac{\mathrm dZ}{\mathrm dx}=X',\quad\frac{\mathrm dZ}{\mathrm du}=U'
\]
Dann gilt identisch
\[
\begin{aligned}
&F(U+wX',\ X-wU')\\
&=F(U,X)\, .\,\psi(U,X,w,L',L'',L'''\text{ usw.})
\end{aligned}
\]
Sooft $U'$ nicht verschwindet, dürfen wir
\[
w=\frac{X-x}{U'}
\]
setzen. Daraus ergibt sich
\[
\begin{aligned}
&F(U+\tfrac{XX'}{U'}-\tfrac{X'x}{U'},\ x)\\
&=F(U,X)\, .\,\psi(U,X,\tfrac{X-x}{U'},L',L'',L'''\text{ usw.})
\end{aligned}
\]
Dies lässt sich auch so aussprechen:

Setzt man in der Funktion $Z$ den Wert $u=U+\frac{XX'}{U'}-\frac{X'x}{U'}$, so geht sie in
\[
F(U,X)\, .\,\psi(U,X,\tfrac{X-x}{U'},L',L'',L'''\text{ usw.})
\]
über.

\SecondArticle{19}
Da im Fall, dass nicht $P=0$ gilt, die Determinante der Funktion $Z$ eine für sich nicht verschwindende Funktion der Unbestimmten $x$ ist, ist die Anzahl der bestimmten Werte von $x$, bei denen diese Determinante den Wert 0 annehmen kann, offenbar endlich. Es lassen sich daher unendlich viele bestimmte Werte von $x$ angeben, die jener Determinante einen von 0 verschiedenen Wert geben. Es sei $X$ ein solcher Wert von $x$, den man außerdem als \textit{reell} annehmen darf. Die Determinante der Funktion $F(u,X)$ ist dann nicht $=0$. Nach Satz II von Artikel 6 können also die Funktionen
\[
F(u,X)\quad\text{und}\quad\frac{\mathrm dF(u,X)}{\mathrm du}
\]
keinen gemeinsamen Teiler haben. Wir nehmen ferner an, es gebe einen bestimmten Wert von $u$, etwa $U$, der entweder reell oder imaginär, also von der Form $g+h\BareRoot{-1}$ ist und $F(u,X)=0$ bewirkt, sodass $F(U,X)=0$ gilt. Dann ist $u-U$ ein Faktor des Polynoms $F(u,X)$, und die Funktion $\frac{\mathrm dF(u,X)}{\mathrm du}$ ist daher gewiss nicht durch $u-U$ teilbar. Nehmen wir also an, diese Funktion,
\SourceLeaf{062}
$\frac{\mathrm dF(u,X)}{\mathrm du}$, erhalte den Wert $U'$, wenn $u=U$ gesetzt wird. Dann kann gewiss nicht $U'=0$ sein. Offenbar ist aber $U'$ der Wert der partiellen Ableitung $\frac{\mathrm dZ}{\mathrm du}$ für $u=U,x=X$. Bezeichnen wir ferner den Wert der partiellen Ableitung $\frac{\mathrm dZ}{\mathrm dx}$ bei denselben Werten von $u,x$ mit $X'$, so ist nach dem im vorigen Artikel Bewiesenen klar, dass die Funktion $Z$ durch die Substitution
\[
u=U+\frac{XX'}{U'}-\frac{X'x}{U'}
\]
identisch verschwindet und daher identisch durch den Faktor
\[
u+\frac{X'}{U'}x-(U+\frac{XX'}{U'})
\]
teilbar ist. Setzen wir nun $u=xx$, so ist offenbar $F(xx,x)$ durch
\[
xx+\frac{X'}{U'}x-(U+\frac{XX'}{U'})
\]
teilbar und nimmt daher den Wert 0 an, wenn für $x$ eine Wurzel der Gleichung
\[
xx+\frac{X'}{U'}x-(U+\frac{XX'}{U'})=0
\]
genommen wird, also wenn man
\[
x=\frac{-X'\pm\BareRoot{(4UU'U'+4XX'U'+X'X')}}{2U'}
\]
setzt. Diese Werte sind bekanntlich entweder reell oder von der Form \mbox{$g+h\BareRoot{-1}$}.

Nun lässt sich leicht beweisen, dass bei denselben Werten von $x$ auch die Funktion $Y$ verschwinden muss. Denn $f(xx,x,\lambda',\lambda'',\lambda'''\text{ usw.})$ ist offenbar das Produkt aus allen $(x-a)(x-b)$ mit ausgeschlossenen Wiederholungen und somit $=\upsilon^{m-1}$. Daraus folgt von selbst
\[
\begin{aligned}
f(xx,x,l',l'',l'''\text{ usw.})&=y^{m-1}\\
f(xx,x,L',L'',L'''\text{ usw.})&=Y^{m-1}
\end{aligned}
\]
oder $F(xx,x)=Y^{m-1}$, dessen bestimmter Wert also nur verschwinden kann, wenn zugleich der Wert von $Y$ verschwindet.

\SecondArticle{20}
Mithilfe der vorangehenden Untersuchungen ist die \textit{Lösung} der Gleichung $Y=0$, also das Auffinden eines bestimmten Wertes von $x$, der entweder reell oder von der Form $g+h\BareRoot{-1}$ ist und ihr genügt, auf die Lösung der Gleichung $F(u,X)=0$ zurückgeführt,
\SourceLeaf{063}
sofern die Determinante der Funktion $Y$ nicht $=0$ ist. Zu bemerken ist: Sind alle Koeffizienten in $Y$, also die Zahlen $L',L'',L'''$ usw., reelle Größen, so sind auch alle Koeffizienten in $F(u,X)$ reell, sofern für $X$ eine reelle Größe genommen wird, was zulässig ist. Der Grad der zweiten Gleichung $F(u,X)=0$ wird durch die Zahl $\frac12m(m-1)$ ausgedrückt. Ist also $m$ eine gerade Zahl von der Form $2^\mu k$, wobei $k$ eine unbestimmte ungerade Zahl bezeichnet, so hat die zweite Gleichung einen Grad von der Form $2^{\mu-1}k$.

Ist die Determinante der Funktion $Y$ $=0$, so lässt sich nach Artikel 10 ein anderes, sie teilendes Polynom $\mathfrak{Y}$ angeben, dessen Determinante nicht $=0$ ist und dessen Grad durch eine Zahl von der Form $2^\nu k$ ausgedrückt wird, wobei entweder $\nu<\mu$ oder $\nu=\mu$ gilt. Jede Lösung der Gleichung $\mathfrak{Y}=0$ genügt auch der Gleichung $Y=0$. Die Lösung der Gleichung $\mathfrak{Y}=0$ wird wiederum auf die Lösung einer anderen Gleichung zurückgeführt, deren Grad durch eine Zahl von der Form $2^{\nu-1}k$ ausgedrückt wird.

Daraus schließen wir also, dass sich allgemein die Lösung jeder Gleichung, deren Grad eine gerade Zahl von der Form $2^\mu k$ ist, auf die Lösung einer anderen Gleichung zurückführen lässt, deren Grad die Form $2^{\mu'}k$ hat, wobei $\mu'<\mu$ gilt. Ist diese Zahl noch immer gerade, also $\mu'$ nicht $=0$, so wird dieselbe Methode erneut angewandt. So fahren wir fort, bis wir zu einer Gleichung gelangen, deren Grad ungerade ist. Alle ihre Koeffizienten sind reell, sofern sämtliche Koeffizienten der ursprünglichen Gleichung reell waren. Eine solche Gleichung ungeraden Grades ist bekanntlich gewiss lösbar, und zwar durch eine reelle Wurzel. Daher sind auch alle vorhergehenden Gleichungen lösbar, entweder durch reelle Wurzeln oder durch Wurzeln von der Form $g+h\BareRoot{-1}$.

\Needspace{16\baselineskip}Damit ist erwiesen, dass jedes Polynom $Y$ von der Form $x^m-L'x^{m-1}+L''x^{m-2}-$ usw., wobei $L',L''$ usw. bestimmte reelle Größen sind, einen Faktor $x-A$ enthält, wobei $A$ eine entweder reelle Größe oder eine Größe von der Form $g+h\BareRoot{-1}$ ist. Im letzteren Fall sieht man leicht, dass $Y$ auch durch die Substitution $x=g-h\BareRoot{-1}$ den Wert 0 erhält, also durch $x-(g-h\BareRoot{-1})$ und daher auch durch das Produkt $xx-2gx+gg+hh$ teilbar ist. Jedes Polynom $Y$ enthält also gewiss einen reellen Faktor ersten oder zweiten Grades. Da dasselbe wiederum für den Quotienten gilt, ist klar, dass sich $Y$ in reelle Faktoren ersten oder zweiten Grades zerlegen lässt. Dies zu beweisen war das Ziel dieser Abhandlung.
\begin{center}\rule{3cm}{0.3pt}\end{center}
